\section{Replication in an independent screen}\label{sec:rep}
This note describes the prespecified replication in the ECCITE-seq screen of Papalexi
et al.\cite{papalexi2021}.

\subsection*{Data, split and protocol}
The scPerturb release\cite{peidli2024} (Zenodo record 7041849) contains 20,729 THP-1 cells stimulated with interferon-$\gamma$,
with knockouts of \RpTargets{} genes and non-targeting guides, three biological
replicates (the 77 cells of a fourth hashtag were not used; \RpCells{} cells
remained), RNA and four surface proteins. Each target-by-replicate group was
split in a fixed pseudo-random order of cells into a training half (RNA and protein
halves) and a test half (adaptation and scoring cells); non-targeting cells were
split into adaptation and scoring cells only. The gene panel was chosen from
training halves by the prespecified rule (\RpGenes{} genes, including the four genes
encoding the antibodies). The protocol and the code were registered after the split and before any statistic of test or
non-targeting cells other than library sizes was computed; the pipeline had been
run on training halves only to check the code. Predictions were computed from
adaptation cells and recorded, and the scoring files were
then opened. There were no amendments.

The prespecified rules of the melanoma screen were reused, with the replicate in the
role of the condition. Every target with an eligible test group was held out in
turn; its channels, controls and paired comparators were fitted on the training
halves of the other targets. For the design criterion, each fold's channel was
also refitted without the five training targets of highest exposure, without the
five of strongest effect and without five random targets (50 draws). The
interferon block comprised \RpIfnGenes{} panel genes of the Hallmark
interferon sets and the antibodies to CD86 and PD-L1.

\subsection*{Results}
Cross-validation over targets chose the smallest penalty scale in every fold,
which retained \RpDimSMedian{} to \RpDimSMax{} supported directions. The five
prespecified criteria were not met: recovery was \RpPrimary\% (95\% CI,
\RpPrimaryLo{} to \RpPrimaryHi\%) over all genes and proteins, lower than both
controls (pseudo-populations \RpAllPseudo\%, derangements \RpAllDerange\%), the
interferon-block share of the paired closed form was \RpIfnShare\%
(\RpIfnShareLo{} to \RpIfnShareHi\%), recovery in non-targeting cells was \RpNtPrimary\%
(\RpNtPrimaryLo{} to \RpNtPrimaryHi\%), and removing the five targets of highest
exposure did not lower recovery relative to random removals (difference
\RpRemDiff\%; \RpRemDiffLo{} to \RpRemDiffHi\%). Paired comparators recovered
\RpPaired\% (closed form) and \RpAllTrans\% (transferred correlation) over all genes
and proteins, and \RpIfnPaired\% and \RpIfnTrans\% in the interferon block. By
replicate, interferon-block recovery of the pairing-free channel was
\RpRepOneIfnPf\%, \RpRepThreeIfnPf\% and \RpRepFourIfnPf\%, against \RpRepOneIfnPaired\%,
\RpRepThreeIfnPaired\% and \RpRepFourIfnPaired\% for the paired closed form.

After the scoring files had been opened, the channel was refitted in every fold with the penalty scale
fixed (exploratory). At 0.1, the value chosen in the
melanoma screen, it retained \RpFixedPointOneDim{} direction and recovered
\RpFixedPointOneAll\% over all genes and proteins and \RpFixedPointOneIfn\% in the
interferon block; at 1, \RpFixedOneDim{} directions, \RpFixedOneAll\% and
\RpFixedOneIfn\%.

\begin{table}[h]
\caption{\textbf{Recovered fraction in the replication screen (\%).} Test targets:
\RpGroups{} target-by-replicate groups from \RpTestTargets{} targets, each held out in
turn. Non-targeting cells: three replicates. Controls and random removals are
averages over their draws.}
\label{tab:rep}
\centering\scriptsize\setlength{\tabcolsep}{3pt}
% Generated by extension/perturbation_replication/make_assets.py. Do not edit.
\begin{tabular}{lccccc}
\toprule
& \multicolumn{3}{c}{Test targets} & \multicolumn{2}{c}{Non-targeting cells} \\
\cmidrule(lr){2-4}\cmidrule(lr){5-6}
Arm & All & Interferon block & Complement & All & Interferon block \\
\midrule
No pairs: perturbations (primary) & \ensuremath{-}212.4 & \ensuremath{-}3.0 & \ensuremath{-}200.5 & \ensuremath{-}89.5 & 20.7 \\
No pairs: corrected RNA & \ensuremath{-}214.4 & \ensuremath{-}6.3 & \ensuremath{-}194.3 & \ensuremath{-}92.7 & 15.4 \\
No pairs: pseudo-populations & \ensuremath{-}10.2 & \ensuremath{-}3.2 & \ensuremath{-}15.6 & \ensuremath{-}15.5 & \ensuremath{-}4.0 \\
No pairs: derangements & \ensuremath{-}45.6 & \ensuremath{-}12.4 & \ensuremath{-}53.6 & \ensuremath{-}25.8 & \ensuremath{-}8.5 \\
No pairs: without 5 highest-exposure targets & \ensuremath{-}172.3 & \ensuremath{-}5.6 & \ensuremath{-}164.0 & \ensuremath{-}132.1 & 13.9 \\
No pairs: without 5 strongest targets & \ensuremath{-}219.3 & \ensuremath{-}22.5 & \ensuremath{-}173.0 & \ensuremath{-}163.7 & \ensuremath{-}7.6 \\
No pairs: without 5 random targets & \ensuremath{-}216.3 & \ensuremath{-}3.7 & \ensuremath{-}226.5 & \ensuremath{-}109.6 & 10.9 \\
Paired closed form & 32.4 & 83.7 & 25.1 & 82.0 & 89.1 \\
Reference regression & 37.1 & 84.3 & 30.6 & 81.3 & 89.4 \\
Transferred correlation & 81.2 & 94.3 & 83.7 & 89.7 & 85.5 \\
Own paired adaptation cells & 21.8 & 60.7 & 3.1 & 74.2 & 71.0 \\
Adaptation half as replicate & \ensuremath{-}204.6 & 18.9 & \ensuremath{-}390.1 & 60.2 & 88.0 \\
\bottomrule
\end{tabular}

\end{table}

\subsection*{Limitations of this test}
The screen had 25 targets, most of them in one pathway, and four proteins, so
the interferon block had only two antibodies. Groups had 40 to 170 adaptation
cells, which limits the within-assay covariance matrices that closed-form transfer
uses.
